% Ad Familiares 4.2 (ad-familiares-04-02) --- English translation
% Translated by Claude (Anthropic) from the Latin (Perseus).
% Style guide: STYLE.md.

\ciceroLetterOpener{M. CICERO S. D. SER. SVLPICIO.}

\ciceroSection{1} On the third day before the Kalends of May, while I was at my house at Cumae, I received your letter; and on reading it I understood that Philotimus had not acted with much sense. Although you had given him instructions, as you write, about every matter, he did not come to me himself but only sent your letter on --- which I gathered had been the shorter for it, since you had supposed that he would carry the rest in person. Still, after I had read your letter, your wife Postumia came to see me, together with our friend Servius. They thought it best that you should come to my place at Cumae, and they also pressed me to write to you to that effect.

\ciceroSection{2} As for the advice you ask of me: it is of the sort I find easier to take myself than to give to another. For what is there I should dare to recommend to a man of the highest authority and the highest prudence? If we are asking what is most upright, the answer is plain; if what is most expedient, obscure. But if we are men of the kind we surely ought to be, who reckon nothing expedient that is not also upright and honourable, then there can be no doubt what we must do.

\ciceroSection{3} As for your sense that my cause is bound up with yours: certainly the error was alike in both of us, and we erred while meaning the very best. For all of our deliberations on either side looked toward concord; and since concord could have been of no greater use to anyone than to Caesar himself, we even reckoned that we would win his favour, remarkably, by defending the peace. How greatly we were deceived, and to what point the situation has now been brought, you see. And you see not only what is being done and what has already been done, but also what the course of events will be and how it will end. Two things therefore remain: either to approve of what is happening, or to take part in it even if you do not approve. Of these the first seems to me shameful, the second downright dangerous.

\ciceroSection{4} What remains is the view that one must withdraw; and within that view there is still room for deliberation --- what plan to follow in withdrawing, what places to make for. Never has the situation been more wretched, and never has deliberation been harder; for nothing can be decided that does not run up against some great difficulty. So my advice, if it seems good to you, is this: if you have already settled in your own mind what you ought to do, and your decision does not depend on mine, then spare yourself the trouble of the journey. But if there is anything you want to talk through with me, I shall wait for you. As far as your own convenience allows, I should wish you to come as soon as may be, since I understood that this is the view of both Servius and Postumia. Farewell.
